\section{Penyederhanaan Ekspresi Logika}

Penyederhanaan ekspresi logika dilakukan dengan menggunakan hukum-hukum logika yang telah dipelajari \cite{rosen2019,libretexts_equiv}. Teknik ini sangat berguna untuk:
\begin{itemize}
  \item Menyederhanakan ekspresi logika yang kompleks menjadi bentuk yang lebih sederhana
  \item Membuktikan ekuivalensi antara dua ekspresi logika
  \item Memeriksa apakah suatu ekspresi merupakan tautologi atau kontradiksi
  \item Mengoptimalkan kondisi dalam program komputer dan rangkaian logika digital
\end{itemize}

\subsection{Teknik Penyederhanaan}

Penyederhanaan dilakukan dengan menerapkan hukum-hukum logika secara berurutan hingga diperoleh bentuk paling sederhana \cite{rosen2019}.

\begin{contoh}
Sederhanakan $(p \wedge \neg q) \vee (p \wedge q \wedge r)$:

\begin{center}
\begin{tabular}{lcl}
  $(p \wedge \neg q) \vee (p \wedge q \wedge r)$ & & \\
  $\equiv$ & $(p \wedge \neg q) \vee (p \wedge (q \wedge r))$ & Tambah Kurung \\
  $\equiv$ & $p \wedge (\neg q \vee (q \wedge r))$ & Distributif \\
  $\equiv$ & $p \wedge ((\neg q \vee q) \wedge (\neg q \vee r))$ & Distributif \\
  $\equiv$ & $p \wedge (T \wedge (\neg q \vee r))$ & Tautologi \\
  $\equiv$ & $p \wedge (\neg q \vee r)$ & Identitas \\
\end{tabular}
\end{center}
\end{contoh}

\subsection{Penyederhanaan untuk Membuktikan Ekuivalensi}

Penyederhanaan juga dapat digunakan untuk membuktikan dua ekspresi logika ekuivalen, yaitu dengan menyederhanakan satu ekspresi hingga menjadi bentuk yang sama dengan ekspresi lainnya \cite{libretexts_equiv,rosen2019}.

\begin{contoh}
Buktikan: $(p \rightarrow q) \wedge (q \rightarrow p) \equiv (p \wedge q) \vee (\neg p \wedge \neg q)$

\begin{center}
\begin{tabular}{lcl}
  $(p \rightarrow q) \wedge (q \rightarrow p)$ & & \\
  $\equiv$ & $(\neg p \vee q) \wedge (\neg q \vee p)$ & $p \rightarrow q \equiv \neg p \vee q$ \\
  $\equiv$ & $(q \vee \neg p) \wedge (p \vee \neg q)$ & Komutatif \\
  $\equiv$ & $(p \vee \neg q) \wedge (q \vee \neg p)$ & Komutatif \\
  $\equiv$ & $((p \vee \neg q) \wedge q) \vee ((p \vee \neg q) \wedge \neg p)$ & Distributif \\
  $\equiv$ & $((p \wedge q) \vee (\neg q \wedge q)) \vee ((p \wedge \neg p) \vee (\neg q \wedge \neg p))$ & Distributif \\
  $\equiv$ & $((p \wedge q) \vee F) \vee (F \vee (\neg q \wedge \neg p))$ & Negasi \\
  $\equiv$ & $(p \wedge q) \vee (\neg q \wedge \neg p)$ & Identitas \\
  $\equiv$ & $(p \wedge q) \vee (\neg p \wedge \neg q)$ & Komutatif \\
\end{tabular}
\end{center}
\end{contoh}

\subsection{Penyederhanaan untuk Memeriksa Tautologi}

Jika penyederhanaan suatu ekspresi logika menghasilkan konstanta $T$ (benar), maka ekspresi tersebut adalah tautologi \cite{rosen2019}.

\begin{contoh}
Periksa apakah $((p \wedge (q \rightarrow r)) \wedge (p \rightarrow (q \rightarrow \neg r))) \rightarrow p$ adalah tautologi:

\begin{center}
\small
\begin{tabular}{lcl}
  & $\neg((p \wedge (\neg q \vee r)) \wedge (\neg p \vee (\neg q \vee \neg r))) \vee p$ & $p \rightarrow q \equiv \neg p \vee q$ \\
  $\equiv$ & $(\neg(p \wedge (\neg q \vee r)) \vee \neg(\neg p \vee (\neg q \vee \neg r))) \vee p$ & De Morgan \\
  $\equiv$ & $((\neg p \vee \neg(\neg q \vee r)) \vee (\neg\neg p \wedge \neg(\neg q \vee \neg r))) \vee p$ & De Morgan \\
  $\equiv$ & $((\neg p \vee (\neg\neg q \wedge \neg r)) \vee (\neg\neg p \wedge (\neg\neg q \wedge \neg\neg r))) \vee p$ & De Morgan \\
  $\equiv$ & $((\neg p \vee (q \wedge \neg r)) \vee (p \wedge (q \wedge r))) \vee p$ & Negasi Ganda \\
  $\equiv$ & $(\neg p \vee (q \wedge \neg r)) \vee (p \wedge (q \wedge r)) \vee p$ & Hapus Kurung \\
  $\equiv$ & $(\neg p \vee (q \wedge \neg r)) \vee ((p \wedge (q \wedge r)) \vee p)$ & Tambah Kurung \\
  $\equiv$ & $(\neg p \vee (q \wedge \neg r)) \vee (p \vee (p \wedge (q \wedge r)))$ & Komutatif \\
  $\equiv$ & $(\neg p \vee (q \wedge \neg r)) \vee p$ & Absorpsi \\
  $\equiv$ & $p \vee (\neg p \vee (q \wedge \neg r))$ & Komutatif \\
  $\equiv$ & $(p \vee \neg p)$ & Absorpsi \\
  $\equiv$ & $T$ & Tautologi \\
\end{tabular}
\end{center}

Karena hasil penyederhanaan adalah $T$, maka ekspresi tersebut adalah \textbf{tautologi}.
\end{contoh}

\begin{catatan}
Penyederhanaan ekspresi logika merupakan keterampilan penting yang memerlukan latihan. Tips:
\begin{enumerate}
  \item Kenali pola-pola umum: $p \wedge \neg p = F$, $p \vee \neg p = T$, $p \wedge T = p$, $p \vee F = p$
  \item Terapkan Hukum De Morgan untuk ``mendorong'' negasi ke dalam
  \item Gunakan distributif untuk mengelompokkan term yang serupa
  \item Carilah pasangan yang menghasilkan $T$ atau $F$ untuk simplifikasi
\end{enumerate}
\end{catatan}
